\chapter{Discussion}
\label{ch:discussion}

This chapter interprets the results, compares the trade-offs between classifiers, acknowledges what this study cannot claim, and discusses the ethical and practical sides of deploying these models.

\section{Interpretation of Results}
\label{sec:interpretation}

After running the experiments, I found that ensemble methods --- Random Forest and XGBoost --- clearly outperformed the three simpler classifiers on both datasets. This makes sense for tabular data with feature interactions: individual trees can learn rules like ``if fever AND rash AND joint pain, predict disease X,'' while a linear model like Logistic Regression cannot represent such conditional patterns without explicit feature engineering.

XGBoost's sequential boosting was particularly helpful for rare diseases. Random Forest can miss rare examples because each tree only sees a bootstrap sample, but XGBoost explicitly focuses each new tree on the mistakes of previous ones. The regularization term in eq.~\ref{eq:xgboost} also helped prevent overfitting on disease categories with few training examples.

All classifiers improved from Dataset~1 to Dataset~2, which confirms what you'd expect from bias-variance theory \citep{Bishop2006}: more data reduces variance. But the benefit was largest for complex models that had room to improve. Naive Bayes hit a ceiling --- its independence assumption limits what it can learn regardless of how much data it sees.

\section{Trade-offs Between Models}
\label{sec:tradeoffs}

Raw performance numbers do not tell the whole story. The choice of classifier for a real deployment depends on several practical factors.

\textbf{Interpretability:} Logistic Regression is the most transparent model. Its coefficients directly show which symptoms push the prediction toward which disease. If a doctor wants to know why the system recommended a dermatologist, Logistic Regression can provide a clear answer. Random Forest and XGBoost offer feature importance rankings but cannot easily explain individual predictions without additional tools like SHAP.

\textbf{Training and prediction speed:} Naive Bayes and Logistic Regression train in seconds even on Dataset~2. SVM was the slowest by a large margin due to its one-vs-one multi-class strategy, which scales poorly with the number of classes. For a mobile health app that needs to update its model frequently, this matters.

\textbf{Handling rare diseases:} XGBoost and Random Forest with class weighting showed the best recall on rare disease classes. Naive Bayes struggled most with imbalanced classes, tending to over-predict common diseases.

\textbf{My recommendation:} For a real healthcare booking system, I would deploy XGBoost or Random Forest as the primary classifier. But I would keep a Logistic Regression model available for when a doctor or regulator asks ``why did the system suggest this specialist?'' --- having an interpretable fallback is worth the slightly lower accuracy.

\section{Limitations and Concerns}
\label{sec:limitations}

This study has several limitations that I want to be upfront about.

\textbf{No external validation.} This is the biggest limitation. I evaluated the classifiers using cross-validation and a held-out test set from the same data distribution. I did not test on an independent clinical dataset. Models that work well on Kaggle data may fail with a different patient population that has different symptom reporting habits or disease prevalence. The accuracy numbers in Chapter~\ref{ch:results} should be treated as optimistic upper bounds.

\textbf{Simplified symptom representation.} The binary encoding doesn't capture severity, duration, or the order symptoms appeared in. A doctor would use all of this. The datasets don't provide it, which limits what the classifiers can learn.

\textbf{Single-disease assumption.} The datasets assign exactly one disease per patient. In reality, patients often have multiple conditions at the same time (comorbidities). A patient with both diabetes and a skin condition would be mishandled by a single-label classifier.

\textbf{No demographic features.} Age, sex, and medical history are important for diagnosis but aren't in either dataset. Adding these would likely improve performance but was outside what the data provided.

\textbf{Hyperparameter search wasn't exhaustive.} I used grid search with predefined ranges. Bayesian optimization or random search over wider ranges might find slightly better configurations, though I doubt this would change the relative ranking of classifiers.

\section{Ethical Considerations}
\label{sec:ethics}

Deploying disease prediction models in healthcare raises issues that go beyond technical performance.

\textbf{Patient safety:} A wrong prediction can lead to a wrong booking. If the model says ``dermatologist'' but the patient actually needs a cardiologist, the delay could be dangerous. Any deployed system must clearly communicate that its predictions are suggestions, not diagnoses, and that seeing a doctor is still necessary.

\textbf{Data protection:} Health data is a special category under the GDPR \citep{GDPR2016}. Processing it requires explicit consent and additional safeguards. A deployed system would need to minimize what data it stores and ensure it is encrypted and access-controlled.

\textbf{EU AI Act:} The recently enacted EU AI Act \citep{EUAI2024} classifies medical AI as high-risk. This means any commercial deployment would require a conformity assessment, risk documentation, human oversight mechanisms, and transparency about how the system works. These are not optional --- they are legal requirements.

\textbf{Bias:} If the training data over-represents certain demographics, the model will perform worse for underrepresented groups. Before deploying this kind of system, fairness audits across age, sex, and ethnicity would be necessary.

\section{Practical Deployment Considerations}
\label{sec:deployment}

Turning these experimental results into a working system involves several steps that are beyond this thesis's scope but worth outlining.

A practical system would present patients with a symptom checklist, encode their selections as a binary vector, run the classifier, and show the predicted disease with a confidence score. The system would then map the disease to the appropriate specialist and suggest available appointment slots.

Not every prediction should be acted on. If the model's confidence is too low, the patient should be directed to a general practitioner instead of a specialist. Setting this threshold means balancing false referrals against missed diagnoses --- a decision that needs medical professionals involved, not just engineers.

Any deployed model would need monitoring. Disease patterns change, new diseases emerge, and the patient population using the system might differ from the training data. Periodic retraining and performance checks would be necessary to catch degradation before it harms patients.
